% !TeX root = ../main.tex
\section*{Общие выводы по работе}
В ходе выполнения практического задания было проведено успешное моделирование устройства цифровой фильтрации и децимации радиосигналов. 
\begin{itemize}
    \item Реализована квадратурная демодуляция, переносящая радиосигнал на нулевую промежуточную частоту с формированием синфазного и квадратурного каналов.
    \item Подтверждено, что динамический диапазон системы (SFDR) ограничивается шумами квантования АЦП, а чрезмерное увеличение разрядности цифрового гетеродина сверх необходимого минимума не дает качественного выигрыша.
    \item Доказана критическая уязвимость архитектуры к I/Q дисбалансу: малейшие фазовые и амплитудные искажения в гетеродине ведут к некомпенсируемому росту зеркальной помехи в полосе пропускания.
    \item Установлено, что замена экономичного CIC-фильтра на качественный КИХ-дециматор позволяет существенно снизить уровень паразитного алиасинга высокочастотных шумов квантования, обеспечив выигрыш в динамическом диапазоне и идеальную плоскость АЧХ при обработке широкополосных ЛЧМ-сигналов.
\end{itemize}